\begin{table}[htbp]\centering
\caption{产品主对比：C3 两种实际可用模式相对原始轨的表现}\label{tab:product}
{\small\begin{threeparttable}
\begin{tabularx}{\linewidth}{Lcc}\toprule
指标 & \makecell{C3 实际 EQ\\（产品预设）} & \textbf{C3 FPS AI} \\
\midrule
脚步中位增益（dB） & $-$0.8 & \hlai{$+$8.7} \\
脚步增益 $\geq$ 3 dB 的比例 & 1\% & \hlai{94\%} \\
枪声中位增益（dB） & $-$1.0 & \hlai{$+$0.2} \\
背景段中位增益（dB） & $-$1.2 & \hlai{$+$0.5} \\
不含脚步的非脚步误增强率 FER & 17\% & \hlai{6\%} \\
枪声频谱总变化 MSE（dB$^{2}$） & 8.3 & \hlai{0.9} \\
枪声频谱形状变化（dB$^{2}$） & 4.4 & \hlai{0.7} \\
含响亮声音片段的最大电平变化（dB） & $-$0.5 & \hlai{$+$0.1} \\
\bottomrule\end{tabularx}
\begin{tablenotes}\footnotesize\item 人工确认事件：脚步 170 个、枪声 72 个、非脚步 87 个；背景段 23 段未经人工核实。增益与频谱变化均相对原始轨计算（定义见表 \ref{tab:metrics}）。\item 加粗：FPS AI 在该项上优于 C3 实际 EQ（脚步增益更高；枪声、背景与响度变化更接近 0；误增强率与频谱变化更低）。\end{tablenotes}\end{threeparttable}}
\end{table}
